%\documentclass[final,3p,times,twocolumn]{elsarticle}
\documentclass[final,3p,times]{elsarticle}
% https://ctan.org/pkg/elsarticle     2018-07-05 
% \documentclass[preprint, review,3p,times,twocolumn]{elsarticle}
% \documentclass[preprint,12pt]{elsarticle}
% Use the option review to obtain double line spacing
% \documentclass[preprint,review,12pt,]{elsarticle}

%\usepackage{paralist, enumerate}
\usepackage[ruled,linesnumbered]{algorithm2e}
\usepackage{amssymb,amsmath,amsthm,bm}
\usepackage{geometry}
\geometry{left=1.8cm,right=1.8cm}

\usepackage{graphicx, subfigure}
%\usepackage{booktabs, multirow, array}
%\usepackage{algorithm, algorithmic}
\usepackage{url}
%\usepackage[switch, mathlines]{lineno}
%\linenumbers\setlength\linenumbersep{1mm}

%\usepackage{enumitem}
%\usepackage{lineno}
%\linenumbers
%\modulolinenumbers[2]

\renewcommand{\vec}[1]{\mathbf{#1}}

%\newtheorem{proposition}{Proposition}
\graphicspath{{figures-pdf/}}
%\newlength\singleimagewidth
%\setlength\singleimagewidth{\columnwidth}

%\usepackage{ulem}

\journal{Optics and Lasers in Engineering}
\usepackage{hyperref}

\begin{document}

\begin{frontmatter}

\title{Cryptanalyzing a bit-level image encryption algorithm based on chaotic maps}

\author[zsi,uestc]{Heping Wen\corref{corr}}
\ead{hepingwen2016@gmail.com}

%\author[uestc]{Chongfu Zhang}

%\author[xtu-cn]{Eric Yong Xie}

\cortext[corr]{Corresponding author.}

\address[zsi]{School of Electronic and Information Engineering, University of Electronic Science and Technology of China, Zhongshan Institute, \\ Zhongshan 528402, China}

\address[uestc]{School of Information and Communication Engineering, University of Electronic Science and Technology of China, Chengdu 611731, China}

\begin{abstract}

Recently, a bit-level image encryption algorithm based on chaotic maps was proposed \cite{LuX2016OLE}, in which a diffusion-permutation structure with only one round is adopted. In this paper, we performed a comprehensive cryptanalysis on the BCIEA. Theoretical analysis and experimental simulations verify that BCIEA cannot resist a chosen ciphertext attack method.


\end{abstract}
\begin{keyword}
image encryption \sep cryptanalysis \sep bit-level \sep chaos.
\end{keyword}
\end{frontmatter}

\section{Introduction}
\label{sec:intro}

%With the continuous rapid development of network communication technology, the security of storage and transmission of multimedia information, especially images and video, is becoming more and more interesting.
%The image has the characteristics of strong correlation and high redundancy of adjacent pixels, thus the traditional cryptographic algorithms, such as AES, DES and IDEA, are not suitable.
%To cope with this challenge, many different methodologies have been studied for exploiting more promising image encryption schemes. 
%However, due to the lack of authoritative metrics, the security of these cryptographic algorithms is still worth exploring and researching.
%
%Nowadays, many image encryption technology has been a hot-spot in the field of information security. In 2019, Ghadirli \textit{et. al} \cite{Ghadirli2019SP} reviewed the current status of color image encryption technology and classified its methodologies into ten categories, including chaos theory, bit plane decomposition, DNA coding, and so on. Yet, comparing with the design of encryption algorithms, the research on cryptanalysis is still relatively lacking. In the same year, Li C. \textit{et. al} \cite{Ghadirli2019SP,cqli:meet:ISAP19} reported the latest image encryption schemes from the viewpoints of a cryptanalyst,and pointed out that it is worth investing more research.
%Our latest works also validates this view \cite{wenhp2019entropy,wenh2019epjp}. In the previous studies, we found that an image encryption scheme \cite{song2015entropy} based on DNA coding can neither resist chosen plaintext attack nor chosen cipher attack, and a bit-level image cipher \cite{Shafique2018epjp} also vulnerable to chosen plaintext attack.
%
%In 2019, Essid \textit{et. al} \cite{ESSAID2019jisa} proposed an image encryption scheme based on a variant of Hill Cipher and 1D chaotic maps, which is abbreviated as BCIEA.
%The autohrs of BCIEA gave some numerical analysis results, including histogram, adjacent pixel correlation, etc., and claimed that it is sufficiently safe.
%However, from the perspective cryptanalysis, several security pitfalls existing in BCIEA can be found as follows:
%\begin{itemize}
%	
%	\item \emph{The existence of an equivalent key}. In BCIEA, the actual secret keys are $(x_0, y_0, z_0, \mu, r_1, r_2)$. Yet, the PRNSs generated by these secret keys are not related to plain images. Accordingly, the PRNSs can be deemed as an equivalent key, exactly $K_1(i),K_2(i),K_3(i)$ for $i =1 \sim L$ and  $h_{11}, h_{12}, h_{21}, h_{22}$.
%	
%	\item \emph{The parameters of the variant Hill cipher is easy to determine}. 
%	In Eqs. (\ref{eq:p1-encryption}) and (\ref{eq:p1-decryption}), considering 
%	In Equation 1, $h_{11},h_{11}^{-1}$ and $(h_{12}K_1(1)-K_2(1))$ can be considered as static unknowns. Based on known/chosen plaintext attacks, the adversary can determine these unknowns by obtaining some pairs of plain and cipher images. Thus, the first pixel can be broken.
%	
%	\item \emph{The keystreams can be determined because of the paradigm of the encryption equations}. Under the premise of obtaining $h_{11}, h_{12}, h_{21}, h_{22}$, the cipher pixels are composed of two separate parts: the plain pixels and the keystreams. Hence, the attackers may solve the keystreams $K_1(i),K_2(i),K_3(i)$ with some pairs of known/chosen plain and cipher images.
%	
%\end{itemize}
%
%Therefore, based on the aforementioned points, we propose an attack method to break BCIEA.

The remainder of the paper is organized as follows. Section~\ref{sec2:BCIEA} reviews the cryptosystem proposed in \cite{ESSAID2019jisa}. Then, Sec.~\ref{sec3:cryptanalysis} presents a comprehensive cryptanalysis on BCIEA. Section~\ref{challenge} verifies the effectiveness of the attack method. The last section concludes the paper.

\section{The cryptosystem under study}
\label{sec2:BCIEA}

This section fist presents the basic principle of bitplane decomposition, and then introduces BCIEA briefly.
\subsection{Bitplane decomposition and composition}


\subsection{Description of BCIEA}

The block diagram of BCIEA is shown in Fig. . As can be seen, the main components of BCIEA are introduced briefly as follows:

%The specific descriptions of CIEA-FOHS are given as follows:
%\begin{figure*}[!htb]
%	\centering
%	\includegraphics[height=4.5cm]{FOHS-flowchart}
%	\caption{The block diagram of BCIEA.}
%	\label{fig:CIEA-FOHS}	
%\end{figure*}   

\begin{itemize}
	
	\item \emph{The chaotic map and the Secret Key}:
	
	In BCIEA, a chaotic map is used twice to generate several keystreams for encryption. The piece-wise linear chaotic map (PWLCM) is mathematically modeled by
	\begin{equation}
	\begin{aligned}
	x_i = F(x_{i-1},\eta)=	\begin{cases}
	x_{i-1}/ \eta,              & 0<x_{i-1}<\eta \\
	(x_{i-1}- \eta)/(0.5-\eta), & \eta \le x_{i-1}<0.5 \\
	F(1-x_{i-1},\eta),          & 0.5 \le x_{i-1}<1 \\
	\end{cases}
	\end{aligned}
	\label{eq:PWLCM}
	\end{equation}
	where $x$ is the state variable, $x_0 \in (0,1)$ is the initial value and $\eta \in (0,0.5)$ is the control parameter, respectively. In BCIEA, two different groups of  initial values and control parameters $(x_0,\eta_1,y_0,\eta_2)$ are served as the secret keys.
	
	\item \emph{Initialization}:
	For the sake of generality, the encrypted object is an 8-bit grayscale image of size $\bm{P}$ of size $M \times N$ (height $\times$ width). The initialization operation consists of two parts: one is converting the plain image into its bit-level form, and the other is generating the keystreams by the chaotic map for the first time.
	\begin{itemize}
		\item By BBD, convert the plain image $\bm{P}$ of size $M \times N$ (height $\times$ width) into two 1D sequences $\bm{A}_1$ and $\bm{A}_2$ of length $4MN$. For convenience, let $L$ be $4MN$.
		\item With the initial parameter $\eta_1$ and the initial value $x_0$, get a sequence $\bm{X}_1$ by iterating Eq. (\ref{eq:PWLCM}) for $N_0+MN$ times and discard the former $N_0$ values, given as
		\begin{equation*}
		X_1(i) = \bmod(floor(x(i))\times 10^{14},256)\\
		\label{eq:PRNS}
		\end{equation*}
		where $i =1 \sim MN$, $floor$ is the function that takes an integer, and $X_1(i) \in \mathbb{Z}_{256}$. Then, two sequences $\bm{b}_1,\bm{b}_2$ of length $L$ are obtained from $\bm{X_1}$ by the method of binary bitplane decomposition. Here, $\bm{b}_1,\bm{b}_2 \in \mathbb{Z}_{256}$.
	\end{itemize}
	
	
	\item \emph{Stage 1. diffusion}:
	
	Step 1. Firstly calculate the sum of $\bm{A}_2$, represented as
	\begin{equation}
	sum_1 = \sum\limits_{i= 1}^{L} A_2(i)
	\label{eq:sum1}
	\end{equation}
	where $L=4MN$. Then run a cyclic right shift operation on $\bm{A}_1$ with $sum_1$ bits to obtain $\bm{A}_{11}$.
	
	Step. 2 Get $\bm{B}_{1}$ by performing the bitwise XOR operations between $\bm{A}_{11}$, $\bm{A}_{2}$ and $\bm{b}_{1}$, given as
	\begin{equation}
	\left\{ \begin{array}{l}
	B_1(1) = A_{11}(1) \oplus A_{11}(L) \oplus A_{2}(1) \oplus b_{1}(1)\\
	B_1(i) = A_{11}(i) \oplus A_{11}(i-1) \oplus A_{2}(i) \oplus b_{1}(i)
	\end{array} \right.
	\label{eq:B1_diffusion}
	\end{equation}
	where $i=2 \sim L$.
	
	Step 3. Like Step 1, firstly compute the sum values $sum_2$ of $\bm{B}_1$ as
	\begin{equation}
	sum_2 = \sum\limits_{i= 1}^{L} B_1(i)
	\label{eq:sum2}
	\end{equation}
	Then , get $\bm{A}_{22}$ cyclic right shifting $\bm{A}_2$ with $sum_2$ bits.
	
	Step. 4 Similar to Step 2, obtain $\bm{B}_{2}$ from $\bm{A}_{22}$, $\bm{B}_{1}$ and $\bm{b}_{2}$, modeled by
	\begin{equation}
	\left\{ \begin{array}{l}
	B_2(1) = A_{22}(1) \oplus A_{22}(L) \oplus B_{1}(1) \oplus b_{2}(1)\\
	B_2(i) = A_{22}(i) \oplus A_{22}(i-1) \oplus B_{1}(i) \oplus b_{2}(i)
	\end{array} \right.
	\label{eq:B2_diffusion}
	\end{equation}
	where $i=2 \sim L$.
	
	\item \emph{Stage 2. Confusion}:
	
	Step 1. Firstly compute the sum values of $\bm{B}_1$ and $\bm{B}_1$ by
	\begin{equation}
	sum = \sum\limits_{i= 1}^{L} (B_1(i)+B_2(i))
	\label{eq:sum3}
	\end{equation}
	Further update the initial value of the chaotic map by
	\begin{equation}
	s_0 = \bmod(y_0+sum/L,1)\\
	\label{eq:update_y0}
	\end{equation}
	Then, with the initial parameter $\eta_2$ and the initial value $s_0$,  using Eq. (\ref{eq:PWLCM}) again, generate a chaotic sequences of length $2L$, and then process it into two integer sequences $\bm{Y}$ and $\bm{Z}$ both of length $L$. Here, $\bm{Y},\bm{Z} \in [1,L]$.
	
	Step 2. Use $\bm{Y}$ and $\bm{Z}$ as coordinate indexes, and then get $\bm{C}_1$ and $\bm{C}_2$ by confusing the positions of the elements in $\bm{B}_2$ and $\bm{B}_1$ respectively. Exactly, given as
	\begin{equation}
	\left\{ \begin{array}{l}
	C_1(i) = B_2(Y(i))\\
	C_2(i) = B_1(Z(i))\\
	\end{array} \right.
	\label{eq:confusion}
	\end{equation}
	where $i=1 \sim L$. 
	
	Finally, get the cipher image $\bm{C}$ from $\bm{C}_1$ and $\bm{C}_2$ by the method of BPC.
\end{itemize}
As stated in \cite{LuX2016OLE}, the decryption process is the inverse of the encryption process, but the specific steps are not described. For more details, please refer to \cite{LuX2016OLE}. 

\section{Cryptanalysis}
\label{sec3:Cryptanalysis}

\subsection{On the decryption of BCIEA}

As everyone knows, ensuring the decryption is available is a basic condition for good cryptosystems. In addition, decryption is an integral part of cryptanalysis. For this reason, we discuss the decryption part firstly.

\begin{itemize}
	
	\item \emph{Initialization}:
	As with encryption, two initialization operations are performed:
	\begin{itemize}
		\item Get $\bm{C}_1$ and $\bm{C}_2$ from the cipher image $\bm{C}$ by BPD. 
		
		\item Generate the keystreams $\bm{b}_1,\bm{b}_2$ by the secret keys $(\eta_1,x_0)$.
	\end{itemize}

	
	
	\item \emph{Stage 1. Anti-confusion}:
	
	Step 1. Firstly compute the sum values of $\bm{C}_1$ and $\bm{C}_2$ by
	\begin{equation}
	sum = \sum\limits_{i= 1}^{L} (C_1(i)+C_2(i))
	\label{eq:sum_C}
	\end{equation}
	Note that the sum of $\bm{B}_1$ and $\bm{B}_2$ is equal to the sum of $\bm{C}_1$ and $\bm{C}_2$ because the confusion only changes the positions of the elements.	Then, with the same secret keys $(\eta_2,y_0)$, update the initial value $y_0$ to $s_0$ and further obtain the two coordinate indexes sequences $\bm{Y}$ and $\bm{Z}$.
	
	Step 2. Restore $\bm{B}_1$ and $\bm{B}_2$ from $\bm{C}_2$ and $\bm{C}_1$ by the anti-confuse, given by
	\begin{equation}
	\left\{ \begin{array}{l}
	B_1(Z(i)) = C_2(i)\\
    B_2(Y(i)) = C_1(i)
	\end{array} \right.
	\label{eq:anti-confusion}
	\end{equation}
	where $i=1 \sim L$. 
	
	
	\item \emph{Stage 2. Anti-diffusion}:
	
	Step 1. Corresponding to Step 4 of the confusion, recover $\bm{A}_{22}$ from $\bm{B}_{2}$, $\bm{B}_{1}$ and $\bm{b}_{2}$, modeled by
	\begin{equation}
	\left\{ \begin{array}{l}
	A_{22}(1) \oplus A_{22}(L)  = B_2(1) \oplus B_{1}(1) \oplus b_{2}(1)\\
	A_{22}(i) \oplus A_{22}(i-1) = B_2(i) \oplus B_{1}(i) \oplus b_{2}(i)
	\end{array} \right.
	\label{eq:B2_to_A22}
	\end{equation}
	where $i=2 \sim L$.
	
	Step 2. Obtain $sum_2$ by Eq. (\ref{eq:sum2}), and then get $\bm{A}_{2}$ by doing cyclic left shift operation on $\bm{A}_22$ with $sum_2$ bits.
	
%	Step 1. Firstly calculate the sum of $\bm{A}_2$, represented as
%	\begin{equation}
%	sum_1 = \sum\limits_{i= 1}^{L} A_2(i)
%	\label{eq:sum1}
%	\end{equation}
%	where $L=4MN$. Then run a cyclic right shift operation on $\bm{A}_1$ with $sum_1$ bits to obtain $\bm{A}_{11}$.
	
	Step. 4 Recover $\bm{A}_{11}$ by performing the bitwise XOR operations between $\bm{B}_{1}$, $\bm{A}_{2}$ and $\bm{b}_{1}$, given as
	\begin{equation}
	\left\{ \begin{array}{l}
	A_{11}(1) \oplus A_{11}(L) = B_1(1) \oplus A_{2}(1) \oplus b_{1}(1)\\
	A_{11}(i) \oplus A_{11}(i-1) B_1(i) \oplus A_{2}(i) \oplus b_{1}(i)
	\end{array} \right.
	\label{eq:B1_to_A11}
	\end{equation}
	where $i=2 \sim L$.
	
	Step 4. Obtain $sum_1$ by Eq. (\ref{eq:sum1}), and then get $\bm{A}_{1}$ by doing cyclic left shift operation on $\bm{A}_11$ with $sum_1$ bits.
	
	Finally, get the cipher image $\bm{C}$ from $\bm{C}_1$ and $\bm{C}_2$ by the method of BPC.
\end{itemize}



%%
%
%\section{Experimental verification}
%\label{sec4:Experiments}
%
%In order to verify the effectiveness of the attack method, some experiments were carried out. The secret key parameters and Image Database selected in the experiments are exactly the same as \cite{ESSAID2019jisa}. Specifically, the six secret keys are: $x_0 = 0.75147854321781$, $y_0 = 0.92791237172793$, $z_0=0.61247892775319$, $\mu=3.98745632101279$, $r_1 = 0.14785214785423$ and $r_2 = 0.61247892775319$. And the USC-SIPI database and the Uncompressed Color Image Database (UCID) are adopted. 
%
%%To test the sensitivity to the encryption key, we encrypt the images “Lena (512  512)”, “Baboon (512 512)” and “Cameramen (256  256)” from the USC-SIPI database [38] , “ucid00416 (512  384)” and “ucid00622 (384 512)” from the UCID (Uncompressed Color Image Database) [56] , with the key:x 0 = 0.75147854321781, 
%The experimental results are shown in Table \ref{tab:exp_results} and Fig. \ref{fig:exp_results}. As can be seen, the proposed attack method can break BCIEA effectively and efficiently.
%
%\begin{figure}[!htb]
%	\centering
%	\subfigure[]{
\includegraphics[height=3cm]{Lena_c}}
%	\subfigure[]{\includegraphics[height=3cm]{Lena}}
%	\subfigure[]{
\includegraphics[height=3cm]{grayCamera_c}}
%	\subfigure[]{\includegraphics[height=3cm]{grayCamera}} \\
%	\subfigure[]{
\includegraphics[height=3cm]{misc512A_c}}
%	\subfigure[]{\includegraphics[height=3cm]{misc512A}}
%	\subfigure[]{
\includegraphics[height=3cm]{misc512B_c}}
%	\subfigure[]{\includegraphics[height=3cm]{misc512B}}
%	\caption{Experimental results for breaking BCIEA.}
%	\label{fig:exp_results}	
%\end{figure}
%
%%\begin{figure}[!htb]
%%	\centering
%%	\subfigure[~Given $\bm{C}$]{\includegraphics[height=\oneimagewid]{misc256A_c}}
%%	\subfigure[~Recovered $\bm{I}$]{\includegraphics[height=\oneimagewid]{misc256A}}
%%	\subfigure[~Given $\bm{C}$]{\includegraphics[height=\oneimagewid]{misc256B_c}}
%%	\subfigure[~Recovered $\bm{I}$]{\includegraphics[height=\oneimagewid]{misc256B}}
%%	\caption{Attacking results with two images of size $256 \times 256 \times 3$ in MISC.}
%%	\label{fig:MISC_256}
%%\end{figure}
%
%\begin{table*}[!htb]
%	\label{tab:exp_result}
%	\caption{Attacking results for breaking LSMCL-IEA.}
%	\centering
%	\begin{tabular}{cccccc}
%		%		\toprule
%		\hline
%		\textbf{Image name} & \textbf{Type and size}  & \textbf{Image database} & \textbf{Encryption time} & \textbf{Breaking time} & \textbf{Data Complexity} \\
%		\hline 
%		Lenna & Gray scale 512 $\times$ 512 & USC-SIPI  & 0.0647s & 0.2539s  & 2  \\ 
%		Cameramen & Gray scale 256 $\times$ 256 & USC-SIPI& 0.0647s & 0.2539s  & 2  \\  
%		Baboon & RGB color 512 $\times$ 512 & USC-SIPI & 0.0647s & 0.2539s  & 2  \\ 
%		ucid00416 & RGB color 512 $\times$ 384 & UCID & 0.0647s & 0.2539s  & 2  \\ 
%		ucid00622 & RGB color 384 $\times$ 512 & UCID & 0.0647s & 0.2539s  & 2  \\ 
%		\hline
%	\end{tabular}
%	\label{tab:exp_results}
%\end{table*}


\section{Conclusion}


This paper performed a comprehensive cryptanalysis on a chaotic image encryption scheme based on a variant of Hill cipher so-called BCIEA. We found that BCIEA cannot resist a chosen plaintext attack because of its inherent security defects. Both theoretical analysis and experimental verification were provided to support that BCIEA can be broken with low computational and data complexity. This study can provide a reference for improving the security of image encryption technology.


\section*{Acknowledgement}

%This research was supported by the Natural Science Foundation of China (No.~61772447, 61532020), the Joint Funds of the National Natural Science Foundation of China, China General Technology Research Institute (No.~U1736113) and the Fundamental Research Funds for the Central Universities (No.~531118010348).

\bibliographystyle{elsarticle-num}
\bibliography{OLE_BIEA}
\end{document}